
où @@M0@@ est un
morphisme lisse. On peut de plus supposer @@M1@@, @@M2@@, @@M3@@ de rang constant, ainsi que les faisceaux conormaux de @@M4@@ et @@M5@@.

Du diagramme commutatif (où tous les morphismes sont lisses) :
@@M6@@
on déduit :
@@M7@@

Soit @@M8@@; alors le rang de @@M9@@ est égal à celui de @@M10@@. Par ailleurs, @@M11@@ définit une section @@M12@@, qui est une immersion régulière de codimension égale au rang de @@M13@@ (EGA IV 17.10.4). Si @@M14@@ est l'immersion déduite de @@M15@@ par le changement de base @@M16@@, on a @@M17@@ et par suite (VII 1.7) :
@@M18@@

En comparant (1.8.1) et (1.8.2), on en déduit la proposition.

\textbf{Définition 1.9.} \uline{L'entier} @@M19@@ \uline{de la proposition} 1.8 \uline{s'appelle} dimension relative virtuelle \uline{de} @@M20@@ \uline{sur} @@M21@@ (\uline{ou de} @@M22@@) \uline{au point} @@M23@@.

La dimension relative virtuelle de @@M24@@ sur @@M25@@ est donc une fonction localement constante sur @@M26@@, à valeurs entières (positives ou négatives). Si @@M27@@ est un morphisme lisse de dimension relative @@M28@@, sa dimension relative virtuelle est @@M29@@. Si @@M30@@ est une immersion régulière de codimension @@M31@@, sa dimension relative virtuelle est @@M32@@.

\textbf{Proposition 1.10.} \uline{Soient} @@M33@@ \uline{et} @@M34@@ \uline{deux morphismes d'intersection complète}, \uline{et soit} @@M35@@ \uline{un point de} @@M36@@. \uline{Alors} :
@@M37@@
